% Page 076

\seite{76}

die Hoffnung aus, daß zukünftig die Fortbildungsschule\ob
in ebenso großer Anzahl und Liebe besucht werde wie\ob
dies bisher der Fall gewesen. An allen Eltern werde er\ob
stets darauf sehen, daß die Kinder die Fortbildungsschule\ob
besuchen und dafür Sorge tragen, daß möglichst alle\ob
zur Abschlußprüfung teilgenommen.


\pend
\begin{verse}
Da der erste Reichspräsident — Ebert — am 28.\ Februar
\end{verse}
\pstart
gestorben ist, wurde eine Reichspräsidentenwahl notwendig.\ob
Der erste Wahltag war der 29.\ März. In unserer\ob
Gemeinde haben von 140 Wahlberechtigten 109 abgestimmt,\ob
2 Stimmen waren ungültig. Von den gültigen\ob
Stimmen entfielen auf Marx (Zentr.) 99 Stimmen,\ob
auf Jarres (Reichsbl.) 5,\ob
auf Held (Bayr. Volkspartei) 2 und auf Braun (Verein\ohb
igte Sozialdemokratische Partei) 1 Stimme.

\margmark{5. April.   D}%
as vorläufige Wahlergebnis der Reichspräsidentenwahl\ob
vom 29.\ 3. lautet folgendermaßen:

\pend
\abschnitt{Braun (Sp.)        7 798 346}
\pstart


\pend
\begin{verse}
Held (Bay. V.)     1 006 790\\
Hellpach (Dem.)    1 567 197\\
Jarres (Reichsbl.)10 408 365\\
Ludendorff (Völk.)   284 975\\
Marx (Z.)          3 884 877\\
Thälmann (Kom.)    1 871 207
\end{verse}
\pstart

Es sind über 26½ Millionen Stimmen abge\ohb
geben worden. Da keiner der Kandidaten die\ob
absolute Majorität erreicht hat, findet am 26.\ April\ob
ein zweiter Wahlgang statt, bei welchem die\ob
relative Majorität entscheiden wird.

Am 31.\ März wurden 11 Schüler entlassen,\ob
darunter 5 Knaben u.\ 6 Mädchen.\ob
Am 1.\ April fand die Neuaufnahme statt.\ob
2 Knaben und 1 Mädchen wurden aufge\ohb
nommen. Somit ist die Schülerzahl von 48 auf\ob
40 zurückgegangen.\ob
Die Osterferien dauerten vom 8.\ bis 20.\ April\ob
einschließlich.
\pend
